% Generated draft: requires mathematical and source review.
\documentclass[11pt]{article}
\usepackage[margin=25mm]{geometry}
\usepackage{fontspec}
\setmainfont{Latin Modern Roman}
\usepackage{amsmath,amssymb,mathtools,unicode-math}
\setmathfont{Latin Modern Math}
\usepackage{enumitem,needspace,xurl,hyperref}
\setlength{\parindent}{0pt}
\setlength{\parskip}{5pt}
\setlength{\emergencystretch}{3em}
\newcommand{\N}{\mathbb{N}}
\newcommand{\Z}{\mathbb{Z}}
\newcommand{\Q}{\mathbb{Q}}
\newcommand{\R}{\mathbb{R}}
\newcommand{\C}{\mathbb{C}}
\newcommand{\F}{\mathbb{F}}
\newcommand{\A}{\mathbb{A}}
\newcommand{\PP}{\mathbb{P}}
\newcommand{\E}{\mathbb{E}}
\newcommand{\eps}{\varepsilon}
\newcommand{\dd}{\,\mathrm{d}}
\begin{document}
\section*{UnsolvedMath ID: 700500}
\textbf{Source problem code:} catalogue entry
\subsection{Definitions and mathematical statement}
For a group \(G\), let \(G'=[G,G]\) and \(G''=[G',G']\), where \([H,H]\) is the subgroup generated by the commutators \(x^{-1}y^{-1}xy\) with \(x,y\in H\). The group \(G\) is metabelian if \(G''=\{1\}\), equivalently if \(G'\) is abelian. It is finitely generated if some finite subset generates it as a group. The integral group ring \(\Z G\) consists of finite formal sums \(\sum_{g\in G}a_g g\), with \(a_g\in\Z\), with multiplication extending the group operation distributively.

Use the convention of unital left \(\Z G\)-modules. A module \(V\) is cyclic if \(V=(\Z G)v\) for some \(v\in V\), and is nonzero if \(V\ne\{0\}\). Equip \(V\oplus V\) with the componentwise \(\Z G\)-action. The supplied source does not specify left or right modules; inversion in \(G\) gives an anti-automorphism of \(\Z G\), so the corresponding existence questions for left and right modules are equivalent.

\textbf{Question (P. M. Neumann).} Does there exist a finitely generated metabelian group \(G\) and a nonzero cyclic \(\Z G\)-module \(V\) such that
\[
V\cong V\oplus V
\qquad\text{as }\Z G\text{-modules}?
\]
Neumann's stated conjecture is that no such pair exists. Equivalently, the conjecture is
\[
\forall G\;\forall V,\qquad
\left.
\begin{gathered}
G\text{ is finitely generated},\qquad G''=\{1\},\\
V\text{ is a nonzero cyclic }\Z G\text{-module}
\end{gathered}
\right\}
\Longrightarrow V\not\cong V\oplus V.
\]

As motivation, the supplied source reports that Neumann exhibited at an Oberwolfach conference in 1979 a finitely generated soluble group of derived length exactly \(3\), together with a nonzero cyclic integral group-ring module isomorphic to its direct sum with itself. Here derived length exactly \(3\) means \(G^{(3)}=\{1\}\) and \(G^{(2)}\ne\{1\}\), with \(G^{(0)}=G\) and \(G^{(n+1)}=[G^{(n)},G^{(n)}]\). This construction is supplied background, not independently verified here. The registry labels the metabelian question open; that label is not an independent assessment of its present status.
\subsection{Short English statement}
Neumann asks whether a finitely generated metabelian group can have a nonzero cyclic integral group-ring module isomorphic, as a module, to two copies of itself. He conjectures that this is impossible. The supplied source reports an example for a finitely generated soluble group of derived length \(3\).
\subsection{Sources}
\begin{itemize}
\item Local UnsolvedMath registry, permanent ID \(700500\), exact source problem code \texttt{catalogue entry}, title \emph{Kourovka Notebook Problem 8.50}. The supplied record attributes the problem to P. M. Neumann, gives the proposal year as 1982, and labels it open. Its recorded update date is \texttt{2026-10-16T14:00:00Z}; this is registry metadata, not a literature-verification date.
\item \emph{The Kourovka Notebook (Unsolved Problems in Group Theory)}, problem 8.50, 8th issue (1982); the supplied background identifies the 21st issue edition (2026), arXiv:\texttt{1401.0300v48}. The registry describes the Issues 1--21 collection as edited by E. I. Khukhro and V. D. Mazurov.
\item P. M. Neumann, contribution in \emph{Groups--Korea, Pusan, 1988}, Lecture Notes in Mathematics, vol. 1398, Springer, Berlin, 1989, pp. 124--139. The problem statement supplies this reference for background and details of the reported construction. No contribution title is supplied.
\end{itemize}
\end{document}
